\documentclass[11pt]{article}

% Encoding and fonts
\usepackage[utf8]{inputenc}
\usepackage[T1]{fontenc}
\usepackage{lmodern}

% Math
\usepackage{amsmath,amssymb,amsthm}
\usepackage{slashed}

% Graphics and colors
\usepackage{graphicx}
\usepackage{xcolor}
\usepackage{caption}
\usepackage{placeins}

% TikZ for diagrams
\usepackage{tikz}
\usetikzlibrary{calc, decorations.pathreplacing, shapes.geometric}

% Better typography
\usepackage{microtype}

% Page layout
\usepackage[left=1in, right=2in, top=1in, bottom=1in, marginparwidth=1.4in, marginparsep=0.2in]{geometry}
\setlength{\headheight}{14pt}
\addtolength{\topmargin}{-2pt}

% Margin notes
\usepackage{marginnote}
\newcommand{\sidefact}[1]{\marginnote{\footnotesize\textcolor{green!50!black}{#1}}}
\newcommand{\sideex}[1]{\marginnote{\footnotesize\textcolor{red!50!black}{#1}}}

% Hyperlinks
\usepackage[hidelinks]{hyperref}

% Lists
\usepackage{enumitem}

% Header: page style
\usepackage{fancyhdr}
\pagestyle{fancy}
\fancyhf{} % clear header and footer
\lhead{\textit{LFT: Lecture 5}}
\rhead{\textit{S. Singh}}
\renewcommand{\headrulewidth}{0.4pt}
\fancyfoot[C]{\thepage} % Add centered page number in footer


% Theorem environments
\newtheorem{theorem}{Theorem}[section]
\newtheorem{lemma}[theorem]{Lemma}
\theoremstyle{definition}
\newtheorem{definition}[theorem]{Definition}

% Metadata
\title{Lecture 6: Some techniques to go around the sign problem}
\author{Simran Singh}
\date{\today}

\begin{document}
\maketitle

\begin{abstract}
In this lecture we revisit and sharpen the distinction between sign quenching
and phase quenching in lattice QCD at finite baryon density, reviewing the
reweighting formulae and the origin of the sign/phase problem in the
thermodynamic limit. We then discuss two conventional strategies for
circumventing the sign problem: the Taylor expansion of thermodynamic
observables about $\mu/T=0$ and the analytic continuation from imaginary
chemical potential. Finally, we introduce \emph{Lefschetz thimbles} as an
alternative approach to the complex-action problem, showing how complexifying the
field variables and integrating along steepest-descent manifolds removes the
oscillatory phase from the path integral.
\end{abstract}


\section{Recap: Sign vs.\ Phase Quenching}


We begin by re-emphasising that \textbf{sign quenching} and \textbf{phase
quenching} are conceptually distinct procedures, even though both are responses
to the complex-action problem that arises in lattice QCD at finite quark
chemical potential $\mu$. The complex-action problem is the statement that for
$\mu\neq 0$ the fermion determinant $\det M(U,\mu)$ is in general complex, so
the integrand $e^{-S_q[U]}\det M(U,\mu)$ cannot be interpreted as a probability
weight for importance sampling. Phase quenching replaces $\det M$ by its
modulus $|\det M|$ (a continuous positive quantity), while sign quenching
replaces it by $\operatorname{Re}\det M$, which is further restricted to positive values by the restriction $|\operatorname{Re}\det M|$ and one keeps only the residual
$\operatorname{sign}$ factor $\pm 1$. The two quenching methods therefore give rise
to genuinely different probability distributions to sample.

%% ──────────────────────────────────────────────────────────────────────────
\subsection{Phase quenching}
%% ──────────────────────────────────────────────────────────────────────────

The full and phase-quenched partition functions are defined respectively as
\begin{align}
  Z_{\text{full}} &= \int \mathcal{D}U\; e^{-S_q[U]}\,\det M(U,\mu) , \\[4pt]
  Z_{\text{pq}}   &= \int \mathcal{D}U\; e^{-S_q[U]}\,\bigl|\det M(U,\mu)\bigr| .
\end{align}
Writing $\det M = |\det M|\,e^{i\phi(U,\mu)}$, the expectation value of an
observable $\mathcal{O}$ in the full theory can be expressed as a ratio of
phase-quenched averages,
\sidefact{The phase $\phi(U,\mu)$ is called the \emph{complex
phase} of the fermion determinant. For $\mu=0$ the $\gamma_5$-Hermiticity of
the Dirac operator forces $\det M$ to be real (in fact non-negative for two
mass-degenerate flavours), so the phase vanishes and there is no sign/phase
problem.}
\begin{equation}
  \langle \mathcal{O}\rangle_{\text{full}}
  = \frac{\langle e^{i\phi}\,\mathcal{O}\rangle_{\text{pq}}}
         {\langle e^{i\phi}\rangle_{\text{pq}}}
  \approx
  \frac{\dfrac{1}{N}\displaystyle\sum_{i=1}^{N}
        e^{i\phi(U_i)}\,\mathcal{O}(U_i)}
       {\dfrac{1}{N}\displaystyle\sum_{i=1}^{N} e^{i\phi(U_i)}} ,
  \label{eq:phase_reweight}
\end{equation}
where the Monte Carlo average on the right is taken with the (positive)
phase-quenched weight $|\det M|\,e^{-S_q}$. The factor $e^{i\phi(U_i)}$ is
evaluated as a post-hoc reweighting on each gauge configuration $U_i$.

\paragraph{What is $\langle e^{i\phi}\rangle_{\text{pq}}$?}

This average, sometimes called the \emph{average phase factor}, controls the
severity of the sign problem. Directly from the definitions above one has
\begin{equation}
  \langle e^{i\phi}\rangle_{\text{pq}}
  = \frac{1}{Z_{\text{pq}}}\int \mathcal{D}U\;e^{-S_q[U]}\,
    |\det M|\,e^{i\phi}
  = \frac{Z_{\text{full}}}{Z_{\text{pq}}} .
  \label{eq:avg_phase}
\end{equation}
Recognising both $Z_{\text{full}}$ and $Z_{\text{pq}}$ as exponentials of
extensive free energies, $Z = e^{-f\,\Omega}$ where $\Omega = V/T$ is the
four-volume, one obtains
\begin{equation}
  \langle e^{i\phi}\rangle_{\text{pq}}
  = e^{-\Delta f\,\Omega} ,
  \qquad \Delta f \equiv f_{\text{full}} - f_{\text{pq}} \geq 0 .
  \label{eq:avg_phase_extensive}
\end{equation}

\paragraph{Why $\Delta f \geq 0$?} By the triangle inequality applied to the path integral,
\begin{equation}
  Z_{\text{full}}
  = \left|\int \mathcal{D}U\;e^{-S_q[U]}\,|\det M|\,e^{i\phi}\right|
  \leq
  \int \mathcal{D}U\;e^{-S_q[U]}\,|\det M|
  = Z_{\text{pq}} .
\end{equation}
Together with $Z = e^{-f\,\Omega}$ this gives
$e^{-f_{\text{full}}\,\Omega} \leq e^{-f_{\text{pq}}\,\Omega}$, which is
equivalent to $f_{\text{full}} \geq f_{\text{pq}}$, i.e.\ $\Delta f \geq 0$.

\paragraph{The sign problem in the thermodynamic limit.}
\sideex{In your exercise you will verify that the relative
statistical error on $\langle\operatorname{sign}\rangle$ in a Monte Carlo
estimate (standard deviation divided by mean) scales as
$e^{\Delta f\,\Omega}/\sqrt{N}$: the
\emph{signal-to-noise ratio} is exponentially suppressed in the volume.
Recovering a fixed precision therefore demands $N \sim e^{2\Delta f\,\Omega}$
samples -- the hallmark of the sign problem.}
Equation~\eqref{eq:avg_phase_extensive} shows that as $\Omega\to\infty$,
\begin{equation}
  \langle e^{i\phi}\rangle_{\text{pq}} = e^{-\Delta f\,\Omega} \to 0 ,
\end{equation}
\textbf{unless} $\Delta f \to 0$, i.e.\ unless the full and phase-quenched
theories happen to be thermodynamically equivalent.  This exponential
suppression of the denominator in
Eq.~\eqref{eq:phase_reweight} is the precise statement of the sign problem
in the phase-quenching approach.

%% ──────────────────────────────────────────────────────────────────────────
\subsection{Sign quenching}
%% ──────────────────────────────────────────────────────────────────────────

In sign quenching one starts from the observation that physical observables
such as $Z_{\text{full}}$ must be real (as argued above). Since the
integration measure is invariant under $U \to U^*$ and the action satisfies
$S_q[U^*] = S_q[U]$, the imaginary part of the integrand cancels pairwise:
\begin{equation}
  Z_{\text{full}}
  = \int \mathcal{D}U\;\det M(U,\mu)\;e^{-S_q[U]}
  = \int \mathcal{D}U\;\operatorname{Re}\!\bigl[\det M(U,\mu)\bigr]\;e^{-S_q[U]} .
  \label{eq:Re_only}
\end{equation}
However, $\operatorname{Re}[\det M]$ is \emph{not} guaranteed to be
non-negative on every configuration -- it merely averages to a positive
number -- so it still cannot serve directly as a probability weight. One
therefore defines the \emph{sign-quenched} partition function by taking the
absolute value pointwise,
\begin{equation}
  Z_{\text{sq}}
  = \int \mathcal{D}U\;\bigl|\operatorname{Re}\!\det M(U,\mu)\bigr|\;e^{-S_q[U]} ,
\end{equation}
\sidefact{ Subscript ``sq'' stands for \emph{sign-quenched}}
and uses the decomposition
\begin{equation}
  \operatorname{Re}\!\det M = \operatorname{sign}(U,\mu)\,
  \bigl|\operatorname{Re}\!\det M\bigr|,
  \qquad \operatorname{sign}(U,\mu) \in \{+1, -1\},
\end{equation}
to write the full expectation value as a reweighted sign-quenched average,
\begin{equation}
  \langle \mathcal{O}\rangle_{\text{full}}
  = \frac{\langle \operatorname{sign}(U)\;\mathcal{O}(U)\rangle_{\text{sq}}}
         {\langle \operatorname{sign}(U)\rangle_{\text{sq}}} .
  \label{eq:sign_reweight}
\end{equation}
For observables that are themselves invariant under $\mathcal{O}(U) =
\mathcal{O}(U^*)$, this representation involves only real numbers and is
well-defined as long as $\langle\operatorname{sign}\rangle_{\text{sq}} \neq
0$. The severity of this condition -- in particular the exponential
suppression $\langle\operatorname{sign}\rangle_{\text{sq}} \sim
e^{-(\Delta f)_{\text{sq}}\,\Omega}$ -- will be examined in the exercises.

\textbf{Phase vs.\ sign quenching, in one line.} Phase quenching uses
$|\det M|$ as the proxy measure and corrects with a continuous phase
$e^{i\phi} \in U(1)$; sign quenching uses $|\!\operatorname{Re}\!\det M|$
and corrects with a discrete sign $\pm 1 \in \mathbb{Z}_2$. Both are
ratio reweightings of the same form
$\langle\mathcal{O}\rangle_{\text{full}} =
\langle w\,\mathcal{O}\rangle_{\text{proxy}}/\langle w\rangle_{\text{proxy}}$.
Moreover, since $|\operatorname{Re}\det M| \leq |\det M|$ pointwise, one has
$Z_{\text{full}} \leq Z_{\text{sq}} \leq Z_{\text{pq}}$ and therefore
\begin{equation*}
  \langle\operatorname{sign}\rangle_{\text{sq}}
  = \frac{Z_{\text{full}}}{Z_{\text{sq}}}
  \;\geq\;
  \frac{Z_{\text{full}}}{Z_{\text{pq}}}
  = \langle e^{i\phi}\rangle_{\text{pq}} :
\end{equation*}
the sign-quenched average reweighting factor is \emph{never smaller} than
the phase-quenched one, i.e.\ $(\Delta f)_{\text{sq}} \leq
(\Delta f)_{\text{pq}}$. 

%% ══════════════════════════════════════════════════════════════════════════
\section{Conventional Approaches at Finite Density}
%% ══════════════════════════════════════════════════════════════════════════

We now survey two established methods for computing observables in lattice
QCD at finite baryon chemical potential $\mu$, both of which avoid simulating
directly at real $\mu$ \cite{Fodor2002, deForcrand2002, DElia2002}.

%% ──────────────────────────────────────────────────────────────────────────
\subsection{Taylor expansion about $\mu/T = 0$}
%% ──────────────────────────────────────────────────────────────────────────
Notice first that the chemical potential always enters as the dimensionless
ratio $\mu/T$. A key exact symmetry of the partition function is
\begin{equation}
  Z(\mu/T) = Z(-\mu/T) ,
  \label{eq:Z_symm}
\end{equation}

which follows from the combined charge-conjugation and $\gamma_5$-Hermiticity
properties of the lattice action. Two exact relations are at work here.
(i) $\gamma_5$-Hermiticity, $\gamma_5 M(U,\mu)\gamma_5 = M(U,-\mu^*)^\dagger$,
gives $\det M[U,\mu] = \det M^*[U,-\mu^*]$ on the \emph{same} configuration,
hence $Z(\mu) = Z^*(-\mu^*)$.
(ii) Charge conjugation, $U_\mu(n)\xrightarrow{C} U_\mu^*(n)$, gives
$\det M[U,\mu] = \det M^*[U^*,\mu^*]$, hence $Z(\mu) = Z^*(\mu^*)$, i.e.\
$Z$ is real for real $\mu$. Combining (i) and (ii) yields
$Z(\mu) = Z(-\mu)$, Eq.~\eqref{eq:Z_symm}.

Since $Z(\mu/T)$ is an even function of $\mu/T$, the pressure
$p = (T/V)\ln Z$ admits the Taylor series \cite{Allton2002}
\begin{equation}
  p = T^4 \sum_{n=0}^{\infty} c_{2n}(T)\left(\frac{\mu}{T}\right)^{2n} ,
  \label{eq:p_taylor}
\end{equation}
or equivalently, writing $\hat p \equiv p/T^4$ and $a_{2n} \equiv c_{2n}(T)$
for brevity,
\begin{align}
  \hat p\!\left(\frac{\mu}{T}\right)
  &= a_0 + a_2\!\left(\frac{\mu}{T}\right)^{\!2}
         + a_4\!\left(\frac{\mu}{T}\right)^{\!4}
         + a_6\!\left(\frac{\mu}{T}\right)^{\!6} + \cdots , \\[4pt]
  \hat p\!\left(\frac{i\mu}{T}\right)
  &= a_0 - a_2\!\left(\frac{\mu}{T}\right)^{\!2}
         + a_4\!\left(\frac{\mu}{T}\right)^{\!4}
         - a_6\!\left(\frac{\mu}{T}\right)^{\!6} + \cdots .
\end{align}
The coefficients $c_{2n}(T)$ are -- in principle -- obtained from
$\mu=0$ simulations via the derivatives of the pressure,
\begin{equation}
  c_{2n}(T) = \frac{1}{(2n)!}\,
  \frac{\partial^{2n}(p/T^4)}{\partial(\mu/T)^{2n}}\bigg|_{\mu=0}.
\end{equation}
This approach is conceptually clean but has two well-known failure modes:
\begin{enumerate}[label=(\arabic*), leftmargin=*, itemsep=4pt]
  \item The higher-order coefficients
        $\partial^{2n}(p/T^4)/\partial(\mu/T)^{2n}\big|_{\mu=0}$ are
        generalised quark-number susceptibilities (cumulants of the
        baryon-number distribution), whose evaluation involves increasingly
        many noisy disconnected diagrams. Their stochastic estimation
        becomes rapidly harder with $n$: beyond $n\!\sim\!4$ the signal
        is typically lost in the statistical noise.
  \item Even if one could compute all coefficients exactly, the series has
        a finite radius of convergence set by the nearest singularity in the
        complex $\mu/T$ plane. One cannot reach physically interesting
        regions of the QCD phase diagram (e.g.\ a putative critical
        end-point) by truncating the series at low order.
\end{enumerate}

%% ──────────────────────────────────────────────────────────────────────────
\subsection{Analytic continuation from imaginary $\mu$}
%% ──────────────────────────────────────────────────────────────────────────
\sidefact{However, at purely imaginary $\mu$, the partition function gets a non-trivial periodicity (Roberge-Weiss symmetry~\cite{Roberge:1986mm}): $Z(\mu_B + i 2 \pi T) = Z(\mu_B)$. This means there is no "new" information outside the period defined by $\mu = 2\pi T$ interval.}

A complementary  strategy exploits the fact that at purely imaginary chemical
potential $\mu = i\mu_I$ with $\mu_I \in \mathbb{R}$, the Dirac operator
satisfies $\gamma_5 M(i\mu_I)\gamma_5 = M(i\mu_I)^\dagger$, so $\det M$
remains real (and positive in the cases of practical interest, e.g. for two degenerate flavours where the measure is $|\det M|^2$, or for staggered fermions where the eigenvalue structure guarantees $\det M > 0$)
\cite{deForcrand2002, DElia2002}. There is therefore no sign problem and one
can perform standard Monte Carlo simulations at $\mu = i\mu_I$. This motivates
performing exact simulations at imaginary $\mu$ and analytically
continuing the results to real $\mu$.

Concretely, one computes the (dimensionless) baryon number density
\begin{equation}
  \frac{n_B}{T^3}\!\left(\frac{\mu}{T},T\right)
  = \frac{\partial\,(p/T^4)}{\partial(\mu/T)}
  = 2a_2\!\left(\frac{\mu}{T}\right)
    + 4a_4\!\left(\frac{\mu}{T}\right)^{\!3}
    + 6a_6\!\left(\frac{\mu}{T}\right)^{\!5} + \cdots ,
\end{equation}
which at imaginary chemical potential takes the form
\begin{equation}
  \frac{n_B}{T^3}\!\left(\frac{i\mu_I}{T},T\right)
  = \frac{i\mu_I}{T}\!\left[2a_2
    - 4a_4\!\left(\frac{\mu_I}{T}\right)^{\!2}
    + 6a_6\!\left(\frac{\mu_I}{T}\right)^{\!4} - \cdots\right] ,
\end{equation}
purely imaginary as required by symmetry. The coefficients $a_{2n}$ can then
be extracted from a fit to lattice data at several values of $\mu_I$ and used
to reconstruct $p(\mu/T)$ at real $\mu$, within the radius of convergence.

%% ══════════════════════════════════════════════════════════════════════════
\section{Lefschetz Thimbles}
%% ══════════════════════════════════════════════════════════════════════════

Both approaches above ultimately rely on avoiding the sign problem
\emph{indirectly}, by simulating at $\mu=0$ or imaginary $\mu$. A more direct
attack on the complex-action problem is offered by the method of
\emph{Lefschetz thimbles} \cite{Cristoforetti2012, Fujii2013}. The idea is
geometric: rather than trying to evaluate an oscillatory integrand on the
original real domain, we \emph{deform} the integration cycle into a complex
manifold on which the imaginary part of the action is constant.

%% ──────────────────────────────────────────────────────────────────────────
\subsection{The oscillation problem}
%% ──────────────────────────────────────────────────────────────────────────
\sidefact{This is a precise version of the sign problem: the
integrand oscillates so rapidly that the Monte Carlo estimate of the integral
has an exponentially small signal-to-noise ratio in the volume.}
Consider the path integral at finite real $\mu$, where the action is complex:
\begin{equation}
  Z(\mu) = \int \mathcal{D}\phi\;
           e^{-S_R[\phi,\mu]\,-\,i S_I[\phi,\mu]} ,
  \qquad S = S_R + i S_I, \;\; S_R, S_I \in \mathbb{R} .
\end{equation}
The fundamental difficulty is that the measure 
$\mathcal{D}\phi = \prod_{i=1}^{N} d\phi_i$ involves $N \sim 10^{6}\text{--}10^{12}$
integrations for a state-of-the-art lattice, and the factor
$e^{-iS_I[\phi,\mu]}$ oscillates wildly in this enormous space, leading to
catastrophic cancellations.
A guiding principle from one-dimensional contour integration is that an
oscillatory integral can often be tamed by deforming the contour into the
complex plane so that the imaginary part of the exponent is stationary --
this is the method of \emph{stationary phase} or \emph{steepest descent}.
The Lefschetz thimble construction is the rigorous generalisation of this
idea to many-dimensional integrals with a holomorphic action.


\textbf{Crucial assumption.} The whole construction requires the action
$S[\phi]$ to be a \emph{holomorphic} function of the complexified fields
$\phi \in \mathbb{C}^N$. For lattice scalar field theory with polynomial
interactions this is automatic. For fermionic theories one integrates out
the fermions first: the resulting \emph{weight} $\det M(U,\mu)\,e^{-S_q}$
is holomorphic in the complexified link variables, since the determinant is
a polynomial in them. Note, however, that the effective \emph{action}
$S_{\text{eff}} = S_q - \ln\det M$ has logarithmic branch points at the
zeros of the determinant; the flow must therefore be defined with the
holomorphic weight, and thimbles can terminate at such zeros.


%% ──────────────────────────────────────────────────────────────────────────
\subsection{Complexification of field variables}
%% ──────────────────────────────────────────────────────────────────────────
\sidefact{The simplest analogy is the passage from a real
line integral $\int_{-\infty}^\infty dx\,f(x)$ to a contour integral
$\int_C dz\,f(z)$ in the complex $z=x+iy$ plane. One looks for a contour
$C$ on which the imaginary part of the exponent is constant -- the
\emph{stationary-phase contour}. Equivalently, one finds a curve
$\beta(x,y) = \text{const}$ where $f(z) = \alpha(x,y) + i\beta(x,y)$.}
By Cauchy's theorem, the integration cycle can be deformed in
\emph{complexified} field space $\phi \in \mathbb{R}^N \hookrightarrow
\mathbb{C}^N$ without changing the value of the integral, provided no
singularities are crossed and the integrand decays at infinity. The question
is: does there exist an alternative integration cycle
\begin{equation}
  \mathcal{D}\phi = \prod_{i=1}^{N}d\phi_i
  \;\longrightarrow\;
  \mathcal{D}\phi_c = \prod_{i=1}^{N}d\phi_{c,i} ,
  \qquad \phi_c = \phi^R + i\phi^I \in \mathbb{C}^N ,
\end{equation}
on which $S_I[\phi_c,\mu]$ is \textbf{constant}? If so, the constant phase
factors out of the integral,
\begin{equation}
  Z(\mu) = e^{-iS_I^{\mathcal{J}}}\int_{\mathcal{J}}
           \mathcal{D}\phi_c\;e^{-S_R[\phi_c,\mu]}
  \quad\Longrightarrow\quad \text{no oscillations in the integrand,}
\end{equation}

where $S_I^{\mathcal{J}}$ is the (constant) value of $S_I$ on the cycle $\mathcal{J}$.
The answer is \textbf{yes}: such cycles exist generically and are called
\emph{Lefschetz thimbles} \cite{Cristoforetti2012}.

%% ──────────────────────────────────────────────────────────────────────────
\subsection{Construction of the thimbles}
%% ──────────────────────────────────────────────────────────────────────────

The construction proceeds in three steps.

\begin{enumerate}[label=(\arabic*), leftmargin=*, itemsep=8pt]

  \item \textbf{Find the critical points.}\;
        The thimbles are attached to the saddle points $\phi_\sigma$ of the
        complexified action, defined by
        \begin{equation}
          \frac{\delta S[\phi_c]}{\delta\phi_c}\Bigg|_{\phi_c=\phi_\sigma}
          = 0 .
        \end{equation}
        

  \item \textbf{Define the steepest-ascent flow.}\;
        \sidefact{\textbf{Sign convention.} We take the
        ascent flow with the $+$ sign so that, by construction, $S_R$
        increases along the trajectory. This is the parametrisation in which
        the thimble is traversed \emph{away} from the saddle as $\tau$
        increases, and the integrand $e^{-S_R}$ is maximal at $\phi_\sigma$.}
        The thimble $\mathcal{J}_\sigma$ associated with the saddle
        $\phi_\sigma$ is the union of all gradient-ascent trajectories of
        $S_R$ starting at $\phi_\sigma$:
        \begin{align}
          \frac{d\phi^R}{d\tau} &= +\frac{\partial S_R}{\partial\phi^R} ,
          \label{eq:flow_R}\\[4pt]
          \frac{d\phi^I}{d\tau} &= +\frac{\partial S_R}{\partial\phi^I} ,
          \label{eq:flow_I}
        \end{align}
        with boundary condition $\phi(\tau\to-\infty) = \phi_\sigma$. This means that the thimble starts at the critical points defined above. The
        \emph{unstable-thimble} $\mathcal{K}_\sigma$ is obtained by reversing
        the sign on the right-hand side: it consists of the
        steepest-\emph{descent} trajectories emanating from the saddle,
        along which $S_R$ decreases.

  \item \textbf{Verify that $S_I$ is constant on $\mathcal{J}_\sigma$.}\;
        Because $S(\phi_c) = S_R + iS_I$ is holomorphic, the Cauchy--Riemann
        equations
        \begin{equation}
          \frac{\partial S_R}{\partial\phi^R}
          = \phantom{-}\frac{\partial S_I}{\partial\phi^I} ,
          \qquad
          \frac{\partial S_R}{\partial\phi^I}
          = -\frac{\partial S_I}{\partial\phi^R}
          \label{eq:CR}
        \end{equation}
        hold pointwise. Along the flow defined by
        Eqs.~\eqref{eq:flow_R}--\eqref{eq:flow_I},
        \begin{align}
          \frac{dS_I}{d\tau}
          &= \frac{\partial S_I}{\partial\phi^R}\frac{d\phi^R}{d\tau}
           + \frac{\partial S_I}{\partial\phi^I}\frac{d\phi^I}{d\tau}
            \notag\\[4pt]
          &= -\frac{\partial S_R}{\partial\phi^I}
             \cdot\frac{\partial S_R}{\partial\phi^R}
           + \phantom{-}\frac{\partial S_R}{\partial\phi^R}
             \cdot\frac{\partial S_R}{\partial\phi^I}
            = 0 ,
        \end{align}
        using \eqref{eq:CR} in the second line. Hence \textbf{$S_I$ is
        constant along every thimble}, as required.

\end{enumerate}


\textbf{Complex form of the flow.} The flow
\eqref{eq:flow_R}--\eqref{eq:flow_I} in real components is equivalent to a
single complex equation in terms of the \emph{Wirtinger derivative}
$\partial/\partial\bar\phi = \tfrac{1}{2}(\partial/\partial\phi^R +
i\,\partial/\partial\phi^I)$. Combining the two real equations via
$d\phi_c/d\tau = d\phi^R/d\tau + i\,d\phi^I/d\tau$ and using the
Cauchy--Riemann relations to write
$\partial S_R/\partial\phi^R + i\,\partial S_R/\partial\phi^I
= \overline{\partial S/\partial\phi_c}$, one obtains the compact form
\begin{equation}
  \frac{d\phi_c}{d\tau}
  = \overline{\left(\frac{\partial S}{\partial\phi_c}\right)} ,
  \qquad
  \frac{d\bar\phi_c}{d\tau}
  = \frac{\partial S}{\partial\phi_c} .
\end{equation}
This is the form most often used in the literature, as it manifests the
holomorphic structure and is well-suited to numerical integration on
$\mathbb{C}^N$.


%% ──────────────────────────────────────────────────────────────────────────
\subsection{Behaviour of $S_R$ along the thimble}
%% ──────────────────────────────────────────────────────────────────────────

Along the ascent flow~\eqref{eq:flow_R}--\eqref{eq:flow_I}, the real part
of the action behaves as
\begin{align}
  \frac{dS_R}{d\tau}
  &= \frac{\partial S_R}{\partial\phi^R}\frac{d\phi^R}{d\tau}
   + \frac{\partial S_R}{\partial\phi^I}\frac{d\phi^I}{d\tau} \notag\\[4pt]
  &= \left(\frac{\partial S_R}{\partial\phi^R}\right)^{\!2}
   + \left(\frac{\partial S_R}{\partial\phi^I}\right)^{\!2}
   = \|\nabla S_R\|^2 \geq 0 .
\end{align}
Hence $S_R$ \emph{increases monotonically} from its value at the saddle
$\phi_\sigma$ as one moves along the thimble. The integrand $e^{-S_R}$ is
therefore \emph{maximal} at the critical point and decays exponentially away
from it, guaranteeing convergence of $\int_{\mathcal{J}_\sigma}
\mathcal{D}\phi_c\,e^{-S_R}$. On the anti-thimble $\mathcal{K}_\sigma$
(steepest descent), $dS_R/d\tau = -\|\nabla S_R\|^2 \le 0$ instead, so $S_R$
\emph{decreases} away from the saddle and the integral diverges -- these
cycles are therefore unsuitable as integration domains.


\textbf{Why this is useful.} On the thimble, $e^{-S_R}$ acts as a genuine
positive weight peaked at the saddle, formally analogous to a Gaussian
centred on $\phi_\sigma$. Importance sampling becomes possible, and the
overall phase $e^{-iS_I^{\mathcal{J}}}$ -- a single constant -- can be pulled out.
The exponentially bad signal-to-noise of the original real-domain integral
is converted into a much milder problem.


%% ──────────────────────────────────────────────────────────────────────────
\subsection{Thimble decomposition and the residual sign problem}
%% ──────────────────────────────────────────────────────────────────────────

In general the original real integration cycle $\mathbb{R}^N$ can be
decomposed as an integer-coefficient sum over stable thimbles attached to
the various saddle points \cite{Cristoforetti2012, Fujii2013, Witten2010},
\sidefact{The integers $m_\sigma$ are intersection numbers
in algebraic topology: $m_\sigma = \langle\mathbb{R}^N,
\mathcal{K}_\sigma\rangle$ counts how many times the original domain
intersects each anti-thimble. They are notoriously hard to compute in
practice and constitute one of the main open problems of the thimble
approach.}
\begin{equation}
  Z(\mu)
  = \sum_{\sigma} m_\sigma\, e^{-iS_{I,\sigma}}\,
    \int_{\mathcal{J}_\sigma} \mathcal{D}\phi_c\;e^{-S_R[\phi_c]} ,
  \qquad m_\sigma \in \mathbb{Z} ,
\end{equation}
so that
\begin{equation}\label{eqn:thimble_decomposition}
  \langle \mathcal{O}\rangle_{\text{full}}
  = \frac{\displaystyle\sum_{\sigma} m_\sigma\, e^{-iS_{I,\sigma}}
          \int_{\mathcal{J}_\sigma} \mathcal{D}\phi_c\;
          e^{-S_R[\phi_c]}\;\mathcal{O}(\phi_c)}
         {\displaystyle\sum_{\sigma} m_\sigma\, e^{-iS_{I,\sigma}}
          \int_{\mathcal{J}_\sigma} \mathcal{D}\phi_c\;e^{-S_R[\phi_c]}} .
\end{equation}
Each individual thimble integral is now \emph{free of oscillations} and the
only remaining phase factors are the global constants $e^{-iS_{I,\sigma}}$.

This rewriting does not, however, eliminate the sign problem entirely -- two
\emph{residual} sources of sign cancellation remain in practice:
\begin{enumerate}[label=(\roman*), leftmargin=*, itemsep=4pt]
  \item \textbf{Inter-thimble interference.} If several thimbles contribute
        with different phases $S_{I,\sigma}$, the sum can suffer cancellations
        analogous to the original sign problem -- though typically much
        milder, since the number of contributing thimbles is small compared
        with the dimension of field space. \textcolor{red}{See Clarifications related to question in lecture at the end of this section.}
  \item \textbf{Jacobian of the parametrisation.} \textcolor{red}{We did not go into this complication in the lecture.} The point is that $\phi \to \phi_c$ is a change of variable -- so there must be an associated Jacobian of the transformation. This can in general lead to a complex phase as well, this is called the residual sign problem. 
\end{enumerate}
Determining \emph{which} thimbles contribute and \emph{with what
multiplicities} $m_\sigma$ (topological intersection numbers) remains very
hard in practice for generic lattice actions, and is the main reason why the
Lefschetz thimble programme has not yet replaced reweighting and analytic
continuation as the workhorse for finite-density lattice QCD.\\

\noindent \textcolor{red}{\textbf{Some clarifications:}}\\
There were some very good questions in the lecture, one of which I only partially answered  and one other question whose answer was only partially correct. Here I offer the correct and hopefully satisfactory answers. Also there was an additional point about the Hessian that you will find in the exercise sheet that I completely forgot to mention:

\begin{itemize}
  \item \textcolor{red}{Hessian:} At a critical point the gradient $\partial S$ vanishes, so the leading non-trivial information about the action's local shape lives in the
\emph{Hessian} $H = \partial^2 S_R$. By holomorphy, $H$ (a real
$2n\times2n$ matrix) has exactly $n$ positive and $n$ negative eigenvalues:
the positive-eigenvalue eigenvectors span the tangent space
$T_{\phi_\sigma}(\mathcal{J}_\sigma)$ -- the $n$ directions of steepest
ascent of $S_R$ \emph{into} the thimble -- while the negative ones span
the anti-thimble. This is why every thimble construction starts at the
saddle: the Hessian gives the only \emph{local} characterisation of which
directions the integration should follow, and it is the seed from which
the rest of the manifold is grown by the flow.
  \item \textcolor{red}{Critical points: } There was a question on the
        importance of the critical points defined by
        $\frac{\delta S[\phi_c]}{\delta\phi_c}\Bigg|_{\phi_c=\phi_\sigma}
        = 0$. This is partially answered in the above point -- the critical
        point is the only place where the gradient vanishes, so the Hessian
        becomes the leading local information about the action and gives us
        a well-defined tangent space to start the integration from. But
        there are deeper reasons too. First, the steepest-descent flow has
        \emph{fixed points} precisely at critical points, so they are the
        only places where flow trajectories can asymptotically end --
        without an endpoint, there is no manifold to build. Second, since
        every trajectory on $\mathcal{J}_\sigma$ terminates at the same
        $\phi_\sigma$, continuity pins $S_I$ to the single value
        $S_I(\phi_\sigma)$ on the entire thimble, which is what allows the
        constant phase $e^{-iS_I}$ to factor out of the integral. Third,
        $S_R$ attains its minimum at $\phi_\sigma$ and grows monotonically
        away from it along $\mathcal{J}_\sigma$, so $e^{-S_R}$ is peaked at
        the saddle and the integral converges automatically. 
 \item \textcolor{red}{Exactness of the decomposition:} Finally, in
        Picard-Lefschetz theory the thimbles attached to critical points
        form a \emph{basis} of the original integration domain, which is what
        allows the decomposition given by Eq.~\eqref{eqn:thimble_decomposition}. It is worth stressing
        that this decomposition is \emph{exact}. This is a direct consequence
        of Cauchy's theorem applied to a closed holomorphic $n$-form, not a
        saddle-point approximation. Restricting to a single thimble
        $\mathcal{J}_0$, by contrast, is \emph{not} exact: it discards every
        other contributing saddle, and its accuracy rests on the conjecture
        that those omitted terms are either non-contributing
        ($n_\sigma = 0$) or exponentially suppressed by
        $e^{-[S_R(\phi_\sigma) - S_R(\phi_0)]}$.

        \textcolor{red}{The correct answer to the question on what makes the $\langle O \rangle_{\mathrm{full}}$ in Eq.~\eqref{eqn:thimble_decomposition} \textbf{real} ?}

        One might worry that each individual term $m_\sigma\,
        e^{-iS_{I,\sigma}}\!\int_{\mathcal{J}_\sigma} \mathcal{D}\phi_c\,
        e^{-S_R}\,\mathcal{O}(\phi_c)$ is complex, so how can the sum
        possibly reproduce a real observable? The resolution is symmetry.
        For real $\mu$, complex conjugation maps saddles to saddles,
        $\phi_\sigma \leftrightarrow \bar\phi_\sigma$, with
        $S(\bar\phi_\sigma) = S(\phi_\sigma)^*$ and
        $\mathcal{J}_{\bar\sigma} = \overline{\mathcal{J}_\sigma}$. Since
        the original cycle $\mathcal{C}$ is itself conjugation-invariant,
        the multiplicities satisfy $m_\sigma = m_{\bar\sigma}$, and the
        contributions pair up:
        \begin{equation*}
          m_\sigma\, e^{-iS_{I,\sigma}}\!\!\int_{\mathcal{J}_\sigma}(\cdots)
          \;+\; m_{\bar\sigma}\, e^{+iS_{I,\sigma}}\!\!
          \int_{\mathcal{J}_{\bar\sigma}}(\cdots)
          \;=\; 2\,\mathrm{Re}\!\left[\,m_\sigma\, e^{-iS_{I,\sigma}}\!\!
          \int_{\mathcal{J}_\sigma}(\cdots)\right] .
        \end{equation*}
        Saddles that are themselves real (fixed under conjugation)
        contribute real amounts on their own. The imaginary parts cancel
        \emph{pairwise and exactly}, just as the imaginary part of $\det M$
        cancels under $U \to U^*$ in the sign-quenching argument. Reality is
        thus an output of symmetry acting on the saddle set, not a property
        of the individual thimble integrals.
\end{itemize}

\noindent \textcolor{red}{\textbf{Lastly}, the conjecture I mentioned in
the lecture about whether the theory is the same after complexification
of the field:} The conjecture is that lattice QCD in the continuum limit
has only one contributing thimble (the one attached to the classical
vacuum). If this is true, all of our problems with complex phases and with
having to account for different thimbles go away.

\medskip
\noindent
But in general: \textit{the original theory is \textbf{not} equivalent to
the new theory restricted to a single thimble. What is exact is the
decomposition given by Eq.~\eqref{eqn:thimble_decomposition}.}
\bibliographystyle{unsrt}
\bibliography{../references}

\end{document}
